% !PRINTED 210
\weberpage{210}
%<<CONT
and from this, by solving, using the notation of \S\, 28, (4):

\weq{5}{m(\mu, \mu') \equiv \begin{pmatrix} a, & b \\ c, & \partial \end{pmatrix} (\nu, \nu').}

Therefore

\weq{6}{\begin{pmatrix} a, & b \\ c, & \partial \end{pmatrix}}

is a suitable symbol for the permutation $S$, and from (5) one sees that two such permutations:

\[
S = \begin{pmatrix} a, & b \\ c, & \partial \end{pmatrix}, \qquad S' = \begin{pmatrix} a', & b' \\ c', & \partial' \end{pmatrix}
\]

compose exactly according to the rule given in \S\, 28:

\weq{7}{S S' = S'' = \begin{pmatrix} a a' + b c', & a b' + b \partial' \\ c a' + \partial c', & c b' + \partial \partial' \end{pmatrix}}

so that the permutation $S''$ of the roots arises when \gesperrt{first} the permutation $S$, and \gesperrt{then} the permutation $S'$, is carried out among the roots of the division equation. It must always be observed, however, that here only the residues of the numbers $a, b, c, \partial$ taken with respect to the modulus $n$ come into consideration, so that the number of permutations $S$ is always finite\footnote{If one wished to choose the notation $S = \begin{pmatrix} \partial, & -b \\ -c, & a \end{pmatrix}$, which at first sight seems to lie nearer, then the composition would have to be denoted by the formula $S'S = S''$, that is, in the reverse sense from the usual composition of permutations % [source typo] German prints "Permentationen".
.}.

1. \gesperrt{The totality of all substitutions $\begin{pmatrix} a, & b \\ c, & \partial \end{pmatrix}$ forms a group, which we denote by $\mathfrak{A}$, and in it, as we have proved, the group $\mathfrak{G}$ of the division equation is contained.}

In $\mathfrak{A}$ there is contained a subgroup $\mathfrak{B}$ which consists of all substitutions (5) % [sic] The German source's cross-reference appears to require (8) or (6).

\weq{8}{T = \begin{pmatrix} \alpha, & \beta \\ \gamma, & \delta \end{pmatrix}}

for which the condition

\weq{9}{\alpha\delta - \beta\gamma \equiv 1 \pmod{n}}

holds.
